% ======================================================
% Compile: xelatex R22_derivation.tex
% Path:    reference/eq17_analysis/derivation/R22_derivation.tex
% Style:   self-contained; mimics Cdpmr_Cn_derivation.tex layout.
% Note:    Companion document: Cdpmr_Cn_derivation.tex (shared notation).
%          S2-S4 (three-layer decomposition, K_var, IF spectral integral /
%          Isserlis) overlap with reference/shared/writeup_architecture.tex
%          S7 (chi_sq, rho_a) and derivation/phase6_R_matrix_derivation.md /
%          phase2_IF_var_dpr_derivation.md.
%          NEW content unique to this document: S5 exact residue/Parseval
%          closed form for rho_dpmr(tau) + S6 3-term geometric IF closed
%          form (verified to <8e-15 against brute force).
% ======================================================
\documentclass[11pt,a4paper]{ctexart}

\usepackage[fleqn]{amsmath}
\usepackage{amssymb,bm}
\setlength{\mathindent}{0pt}
\usepackage{geometry}
\geometry{margin=2.2cm}
\usepackage{enumitem}
\usepackage{array,booktabs}
\usepackage{titlesec}
\usepackage{hyperref}
\hypersetup{colorlinks=true,linkcolor=blue,citecolor=blue}

\setlength{\parindent}{0pt}
\setlength{\parskip}{0.15em}
\setlength{\abovedisplayskip}{4pt}
\setlength{\belowdisplayskip}{4pt}
\setlength{\abovedisplayshortskip}{2pt}
\setlength{\belowdisplayshortskip}{2pt}
\titlespacing*{\section}{0pt}{1.4ex}{0.6ex}
\titlespacing*{\subsection}{0pt}{0.8ex}{0.3ex}

% --- Inherited symbols ---
\newcommand{\lc}{\lambda_c}
\newcommand{\dxT}{\delta x_{T}}
\newcommand{\nx}{n_{x}}
\newcommand{\dpmr}{\delta p_{mr}}
\newcommand{\FT}{F_{T}}
\newcommand{\FN}{F_{N}}
\newcommand{\HE}{H_{E}}
\newcommand{\Cdpmr}{C_{\mathrm{dpmr}}}
\newcommand{\Cn}{C_{n}}
\newcommand{\apd}{a_{pd}}
\newcommand{\acov}{a_{cov}}
\newcommand{\sighatsq}{\hat{\sigma}^{2}}

% --- New symbols ---
\newcommand{\Kvar}{K_{\mathrm{var}}}
\newcommand{\ggauss}{g_{\mathrm{gauss}}}
\newcommand{\gE}{g_{E}}
\newcommand{\IF}{\mathrm{IF}}
\newcommand{\Rdpmr}{R_{\dpmr}}
\newcommand{\rhodpmr}{\rho_{\dpmr}}
\newcommand{\RfT}{R_{\FT}}
\newcommand{\RfN}{R_{\FN}}

\begin{document}
\thispagestyle{empty}

\begin{center}
{\Large\bfseries Derivation of $R(2,2)$}
\end{center}
\vspace{0.3cm}

%=====================================================================
\section{Setup}
%=====================================================================

\subsection{HP Residual and Transfer Functions}

For closed-loop delay $d=2$, the HP residual is
\begin{equation}
\dpmr(z^{-1}) = -\FT(z^{-1})\,\dxT(z^{-1}) + \FN(z^{-1})\,\nx(z^{-1}),
\label{eq:dpmr}
\end{equation}
\begin{align}
\FT(z^{-1}) &= \frac{(1-\apd)(1-z^{-1})\,z^{-3}\,[\,1+(1-\lc)z^{-1}+(1-\lc)z^{-2}\,]}{(1-(1-\apd)z^{-1})(1-\lc z^{-1})}, \label{eq:FT}\\[2pt]
\FN(z^{-1}) &= (1-z^{-1})\,z^{3}\,\FT(z^{-1}), \label{eq:FN-vs-FT}
\end{align}
where $\dxT,\nx$ are independent zero-mean Gaussian white noise with variances $\sigma^{2}_{\dxT},\sigma^{2}_{\nx}$; $\apd\in(0,1)$ is the IIR LP coefficient and $\lc\in(0,1)$ is the closed-loop pole.

\subsection{Variance of $\dpmr$}

From the companion document,
\begin{equation}
\sigma^{2}_{\dpmr} = \Cdpmr(\lc,\apd)\,\sigma^{2}_{\dxT} + \Cn(\lc,\apd)\,\sigma^{2}_{\nx}.
\label{eq:sigma-dpmr}
\end{equation}

\subsection{Variance Estimator and Target}

\begin{equation}
\sighatsq[k] = (1-\acov)\,\sighatsq[k-1] + \acov\,\dpmr^{2}[k],
\qquad
\HE(z^{-1}) = \frac{\acov}{1-(1-\acov)z^{-1}},\quad \acov\in(0,1).
\label{eq:EWMA}
\end{equation}
\begin{equation}
R(2,2) := \mathrm{Var}\bigl(\sighatsq[k]\bigr)\quad\text{at }k\to\infty.
\label{eq:R22-def}
\end{equation}

%=====================================================================
\section{Three-Layer Decomposition}
%=====================================================================

\begin{equation}
R(2,2) = \Kvar\,\cdot\,\sigma^{4}_{\dpmr}\,\cdot\,\IF,
\label{eq:R22-decomp}
\end{equation}
\begin{itemize}[leftmargin=*,topsep=0pt,itemsep=0pt]
\item $\Kvar$ : variance-estimation prefactor, function of $\acov$ only (\S\ref{sec:Kvar}).
\item $\sigma^{2}_{\dpmr}$ : from~\eqref{eq:sigma-dpmr}.
\item $\IF$ : color-inflation factor ($\IF=1$ for white $\dpmr$), \S\S\ref{sec:IF-int}--\ref{sec:IF-closed}.
\end{itemize}

%=====================================================================
\section{Closed Form for $\Kvar$}
\label{sec:Kvar}
%=====================================================================

\subsection{Gaussian Fourth-Moment Factor}

By Isserlis' theorem, for $X\sim\mathcal{N}(0,\sigma^{2})$, $\mathrm{Var}(X^{2}) = 2\sigma^{4}$. Define
\begin{equation}
\ggauss := \mathrm{Var}(\dpmr^{2})\,/\,\sigma^{4}_{\dpmr} = 2.
\label{eq:ggauss}
\end{equation}

\subsection{Variance Gain of $\HE$}

For a stable LTI filter $H$ and stationary input $x$ with PSD $S_{x}$,
\begin{equation}
\mathrm{Var}(H\cdot x) = \frac{1}{2\pi}\!\!\int_{-\pi}^{\pi}\!\!|H(e^{j\omega})|^{2}\,S_{x}(\omega)\,d\omega,
\label{eq:univ-var}
\end{equation}
which for white $x$ reduces to $\|H\|^{2}\sigma^{2}_{x}$ with variance gain $\|H\|^{2}:=(1/2\pi)\int|H|^{2}d\omega$. Applied to $\HE$ in~\eqref{eq:EWMA},
\begin{equation}
\gE := \|\HE\|^{2}
     = \frac{1}{2\pi}\!\!\int_{-\pi}^{\pi}\!\!|\HE(e^{j\omega})|^{2}\,d\omega
     = \frac{\acov}{2-\acov}.
\label{eq:gE}
\end{equation}

\subsection{Composition}
\begin{equation}
\Kvar := \ggauss\cdot\gE = \frac{2\,\acov}{2-\acov}.
\label{eq:Kvar}
\end{equation}

%=====================================================================
\section{$\IF$ as a Spectral Integral}
\label{sec:IF-int}
%=====================================================================

\subsection{Apply Universal Variance Formula to $\sighatsq$}

Let $v[k] := \dpmr^{2}[k] - \sigma^{2}_{\dpmr}$ be the centered squared residual, with PSD $S_{v}(\omega)$. Applying~\eqref{eq:univ-var} with $H=\HE$,
\begin{equation}
R(2,2) = \frac{1}{2\pi}\!\!\int_{-\pi}^{\pi}\!\!|\HE(e^{j\omega})|^{2}\,S_{v}(\omega)\,d\omega.
\label{eq:R22-integral}
\end{equation}
Because $\dpmr$ inherits the LTI poles $1-\apd$ and $\lc$ from $\FT,\FN$, $v$ is colored ($S_{v}$ non-constant).

\subsection{Isserlis' Theorem at Lag $\tau$}

Define
\begin{equation}
\Rdpmr(\tau) := \mathrm{Cov}\bigl(\dpmr[k],\dpmr[k+\tau]\bigr),
\qquad
\rhodpmr(\tau) := \Rdpmr(\tau)/\sigma^{2}_{\dpmr}.
\label{eq:Rdpmr-def}
\end{equation}
By Isserlis' theorem applied at lag $\tau$ for zero-mean stationary Gaussian $\dpmr$,
\begin{equation}
R_{v}(\tau) = 2\,\Rdpmr^{2}(\tau),\qquad
\rho_{v}(\tau) = \rhodpmr^{2}(\tau).
\label{eq:isserlis}
\end{equation}

\subsection{Result}

Evaluating the integral~\eqref{eq:R22-integral} using \eqref{eq:isserlis} (details omitted),
\begin{equation}
\boxed{\;
R(2,2) = \Kvar\cdot\sigma^{4}_{\dpmr}\cdot\IF,
\qquad
\IF := 1 + 2\!\sum_{\tau\ge 1}\rhodpmr^{2}(\tau)\,s^{\tau},
\qquad s := 1-\acov.
\;}
\label{eq:IF-def}
\end{equation}
$\IF=1$ for white $\dpmr$ ($\rhodpmr(\tau)=0$, $\tau\ge 1$); the colored case is closed in \S\S\ref{sec:rho}--\ref{sec:IF-closed}.

%=====================================================================
\section{Closed Form for $\rhodpmr(\tau)$}
\label{sec:rho}
%=====================================================================

\subsection{Local Notation}
\begin{equation}
\alpha := 1-\apd,\qquad \beta := \lc,\qquad c := 1-\lc.
\label{eq:alpha-beta}
\end{equation}

\subsection{$\RfT(\tau)$ via Parseval and Residues}

By Parseval,
\begin{equation}
\RfT(\tau) = \frac{1}{2\pi}\!\!\int_{-\pi}^{\pi}\!\!|\FT(e^{j\omega})|^{2}\cos(\omega\tau)\,d\omega.
\label{eq:RfT-int}
\end{equation}
Evaluating~\eqref{eq:RfT-int} by residues at the poles $z=\alpha$ and $z=\beta$ of $\FT$,
\begin{equation}
\RfT(\tau) = A_{\alpha}\,\alpha^{\tau} + A_{\beta}\,\beta^{\tau},\qquad \tau\ge 3,
\label{eq:RfT-tail}
\end{equation}
with $A_{\alpha},A_{\beta}$ closed in $(\alpha,\beta,c)$. Values $\RfT(0),\RfT(1),\RfT(2)$ follow by direct evaluation of~\eqref{eq:RfT-int}.

\subsection{Autocorrelation $\RfN(\tau)$}

From~\eqref{eq:FN-vs-FT}, $|\FN(e^{j\omega})|^{2} = 2(1-\cos\omega)|\FT(e^{j\omega})|^{2}$. By inverse DTFT,
\begin{equation}
\RfN(\tau) = 2\,\RfT(\tau) - \RfT(\tau-1) - \RfT(\tau+1).
\label{eq:RfN-id}
\end{equation}

\subsection{Combined $\rhodpmr(\tau)$}

By independence of $\dxT,\nx$,
\begin{equation}
\Rdpmr(\tau) = \sigma^{2}_{\dxT}\,\RfT(\tau) + \sigma^{2}_{\nx}\,\RfN(\tau),
\qquad
\rhodpmr(\tau) = \Rdpmr(\tau)/\sigma^{2}_{\dpmr}.
\label{eq:rho-comb}
\end{equation}
For $\tau\ge 3$, both $\RfT,\RfN$ are linear combinations of $\alpha^{\tau},\beta^{\tau}$, so
\begin{equation}
\rhodpmr(\tau) = c_{\alpha}\,\alpha^{\tau} + c_{\beta}\,\beta^{\tau},\qquad \tau\ge 3,
\label{eq:rho-tail}
\end{equation}
with $c_{\alpha},c_{\beta}$ closed in $(\alpha,\beta,c,\sigma^{2}_{\dxT},\sigma^{2}_{\nx})$. Values $\rhodpmr(1),\rhodpmr(2)$ follow by case-split evaluation.

%=====================================================================
\section{Closed Form for $\IF$}
\label{sec:IF-closed}
%=====================================================================

Squaring~\eqref{eq:rho-tail},
\begin{equation}
\rhodpmr^{2}(\tau) = c_{\alpha}^{2}\,(\alpha^{2})^{\tau} + 2c_{\alpha}c_{\beta}\,(\alpha\beta)^{\tau} + c_{\beta}^{2}\,(\beta^{2})^{\tau},\qquad \tau\ge 3.
\label{eq:rho-sq}
\end{equation}
The $\tau\ge 3$ part of the sum in~\eqref{eq:IF-def} splits into three geometric series:
\begin{equation}
\sum_{\tau\ge 3}(\alpha^{2}s)^{\tau} = \frac{(\alpha^{2}s)^{3}}{1-\alpha^{2}s},
\end{equation}
and likewise for $\alpha\beta s,\beta^{2}s$. With the explicit $\tau=1,2$ terms,
\begin{equation}
\boxed{\;
\IF = 1 + 2\!\left[\,
\rhodpmr^{2}(1)\,s + \rhodpmr^{2}(2)\,s^{2}
+ c_{\alpha}^{2}\frac{(\alpha^{2}s)^{3}}{1-\alpha^{2}s}
+ 2c_{\alpha}c_{\beta}\frac{(\alpha\beta s)^{3}}{1-\alpha\beta s}
+ c_{\beta}^{2}\frac{(\beta^{2}s)^{3}}{1-\beta^{2}s}
\,\right].
\;}
\label{eq:IF-closed}
\end{equation}
Convergence requires $\alpha^{2}s,\alpha\beta s,\beta^{2}s<1$, guaranteed by $\alpha,\beta,s\in(0,1)$.

\textit{Caveat.} The form~\eqref{eq:rho-tail} is valid only for $\tau\ge 3$. Replacing the three $(\cdot)^{3}/(1-\cdot)$ tails by $(\cdot)/(1-\cdot)$ and dropping the explicit $\tau=1,2$ terms introduces relative errors up to $\sim 135\%$ at $\apd=0.5$.

\end{document}
